\textbf{Phần 1. Trắc nghiệm lựa chọn}
\setcounterpageref{ex}{0}

\begin{ex} %[5D6V31]
Giá trị của biểu thức $4{,}75 \times 3{,}8 + 4{,}75 \times 6{,}2$ là
\choice
{$4{,}75$}
{\True $47{,}5$}
{$475$}
{$48$}
\loigiai{
Đặt thừa số chung $4{,}75$:
$$4{,}75 \times (3{,}8 + 6{,}2) = 4{,}75 \times 10 = 47{,}5.$$
\textbf{Đáp án của câu hỏi là B}
}
\end{ex}

\begin{ex} %[5D6V31]
Tìm $x$, biết: $x \times 2{,}4 + x \times 7{,}6 = 68$.
\choice
{$0{,}68$}
{\True $6{,}8$}
{$68$}
{$7{,}2$}
\loigiai{
Ta có:
\begin{align*}
x \times (2{,}4 + 7{,}6) &= 68\\
x \times 10 &= 68\\
x &= 68 : 10 = 6{,}8
\end{align*}
\textbf{Đáp án của câu hỏi là B}
}
\end{ex}

\begin{ex} %[5D6V31]
Tính giá trị của biểu thức: $(12{,}5 - 4{,}25) : 1{,}5 + 4{,}5$.
\choice
{$8$}
{$9$}
{\True $10$}
{$11$}
\loigiai{
Ta có:
$$(12{,}5 - 4{,}25) : 1{,}5 + 4{,}5 = 8{,}25 : 1{,}5 + 4{,}5 = 5{,}5 + 4{,}5 = 10.$$
\textbf{Đáp án của câu hỏi là C}
}
\end{ex}

\begin{ex} %[5D6V31]
Một cửa hàng ngày thứ nhất bán được $125{,}5\text{ kg}$ gạo. Ngày thứ hai bán nhiều hơn ngày thứ nhất $24{,}8\text{ kg}$. Ngày thứ ba bán ít hơn ngày thứ hai $15{,}3\text{ kg}$. Cả ba ngày cửa hàng bán được bao nhiêu ki-lô-gam gạo?
\choice
{$400{,}8\text{ kg}$}
{\True $410{,}8\text{ kg}$}
{$412{,}5\text{ kg}$}
{$395{,}6\text{ kg}$}
\loigiai{
Ngày thứ hai bán được:
$125{,}5 + 24{,}8 = 150{,}3\text{ (kg)}$.\\
Ngày thứ ba bán được:
$150{,}3 - 15{,}3 = 135\text{ (kg)}$.\\
Cả ba ngày bán được:
$125{,}5 + 150{,}3 + 135 = 410{,}8\text{ (kg)}$.\\
\textbf{Đáp án của câu hỏi là B}
}
\end{ex}

\begin{ex} %[5D6V31]
Trung bình cộng của 4 số là $14{,}5$. Biết tổng của ba số đầu là $42{,}3$. Số thứ tư là
\choice
{$14{,}7$}
{\True $15{,}7$}
{$16{,}7$}
{$15{,}2$}
\loigiai{
Tổng của cả 4 số là:
$14{,}5 \times 4 = 58$.\\
Số thứ tư là:
$58 - 42{,}3 = 15{,}7$.\\
\textbf{Đáp án của câu hỏi là B}
}
\end{ex}

\begin{ex} %[5D6V31]
Khi nhân một số với $2{,}5$, một học sinh sơ ý chia số đó cho $2{,}5$ nên nhận được kết quả là $1{,}6$. Kết quả đúng của phép nhân ban đầu là
\choice
{$4$}
{$8$}
{\True $10$}
{$12{,}5$}
\loigiai{
Số đem nhân ban đầu là:
$1{,}6 \times 2{,}5 = 4$.\\
Kết quả đúng của phép nhân là:
$4 \times 2{,}5 = 10$.\\
\textbf{Đáp án của câu hỏi là C}
}
\end{ex}

\begin{ex} %[5D6V31]
Mua 5 quyển vở và 3 chiếc bút cùng loại hết $64\,000$ đồng. Mua 5 quyển vở và 5 chiếc bút cùng loại hết $80\,000$ đồng. Giá tiền của một quyển vở là
\choice
{$6\,000\text{ đồng}$}
{\True $8\,000\text{ đồng}$}
{$10\,000\text{ đồng}$}
{$12\,000\text{ đồng}$}
\loigiai{
Giá tiền của 2 chiếc bút là:
$80\,000 - 64\,000 = 16\,000\text{ (đồng)}$.\\
Giá tiền của 1 chiếc bút là:
$16\,000 : 2 = 8\,000\text{ (đồng)}$.\\
Giá tiền của 5 quyển vở là:
$64\,000 - 8\,000 \times 3 = 40\,000\text{ (đồng)}$.\\
Giá tiền của 1 quyển vở là:
$40\,000 : 5 = 8\,000\text{ (đồng)}$.\\
\textbf{Đáp án của câu hỏi là B}
}
\end{ex}

\begin{ex} %[5D6V31]
Tính thuận tiện: $0{,}25 \times 12{,}4 \times 40$.
\choice
{$12{,}4$}
{\True $124$}
{$1\,240$}
{$31$}
\loigiai{
Ta có:
$$0{,}25 \times 12{,}4 \times 40 = (0{,}25 \times 40) \times 12{,}4 = 10 \times 12{,}4 = 124.$$
\textbf{Đáp án của câu hỏi là B}
}
\end{ex}

\begin{ex} %[5D6V31]
Một ô tô đi quãng đường $120\text{ km}$ hết $2{,}5$ giờ. Với vận tốc đó, để đi hết quãng đường $168\text{ km}$ thì ô tô mất thời gian là
\choice
{$3\text{ giờ}$}
{\True $3{,}5\text{ giờ}$}
{$3{,}8\text{ giờ}$}
{$4\text{ giờ}$}
\loigiai{
Vận tốc của ô tô là:
$120 : 2{,}5 = 48\text{ (km/h)}$.\\
Thời gian để ô tô đi hết $168\text{ km}$ là:
$168 : 48 = 3{,}5\text{ (giờ)}$.\\
\textbf{Đáp án của câu hỏi là B}
}
\end{ex}

\begin{ex} %[5D6V31]
Tìm số dư của phép chia $123{,}4 : 5{,}6$ nếu phần thập phân của thương lấy đến 2 chữ số.
\choice
{$0{,}32$}
{$3{,}2$}
{\True $0{,}032$}
{$0{,}0032$}
\loigiai{
Thực hiện phép chia: $123{,}4 : 5{,}6 = 22{,}03$ (lấy đến 2 chữ số thập phân).\\
Ta có: $22{,}03 \times 5{,}6 = 123{,}368$.\\
Số dư là:
$123{,}4 - 123{,}368 = 0{,}032$.\\
\textbf{Đáp án của câu hỏi là C}
}
\end{ex}

\begin{ex} %[5D6V31]
Bác An mua $2{,}5\text{ kg}$ thịt lợn hết $275\,000$ đồng và $1{,}5\text{ kg}$ thịt bò hết $360\,000$ đồng. Hỏi mỗi ki-lô-gam thịt bò đắt hơn mỗi ki-lô-gam thịt lợn bao nhiêu tiền?
\choice
{$110\,000\text{ đồng}$}
{$120\,000\text{ đồng}$}
{\True $130\,000\text{ đồng}$}
{$140\,000\text{ đồng}$}
\loigiai{
Giá 1 kg thịt lợn là:
$275\,000 : 2{,}5 = 110\,000\text{ (đồng)}$.\\
Giá 1 kg thịt bò là:
$360\,000 : 1{,}5 = 240\,000\text{ (đồng)}$.\\
Mỗi ki-lô-gam thịt bò đắt hơn thịt lợn là:
$240\,000 - 110\,000 = 130\,000\text{ (đồng)}$.\\
\textbf{Đáp án của câu hỏi là C}
}
\end{ex}

\begin{ex} %[5D6V31]
Tính giá trị của biểu thức: $100 - 24{,}6 : 0{,}6 + 5{,}8$.
\choice
{$54{,}8$}
{$63{,}8$}
{\True $64{,}8$}
{$65{,}8$}
\loigiai{
Ta có:
$$100 - 24{,}6 : 0{,}6 + 5{,}8 = 100 - 41 + 5{,}8 = 59 + 5{,}8 = 64{,}8.$$
\textbf{Đáp án của câu hỏi là C}
}
\end{ex}

\textbf{Phần 2. Trắc nghiệm đúng sai}
\setcounterpageref{ex}{0}

\begin{ex} %[5D6V31]
Cho hai biểu thức: $X = 3{,}6 \times 4{,}5 + 3{,}6 \times 5{,}5$ và $Y = 72 : 1{,}8 - 4$.
\choiceTF
{\True Giá trị của biểu thức $X$ bằng $36$}
{\True Giá trị của biểu thức $Y$ bằng $36$}
{\True Hai biểu thức $X$ và $Y$ có giá trị bằng nhau}
{\True Tích $X \times Y$ bằng $1\,296$}
\loigiai{
\begin{itemchoice}
\item Đúng. $X = 3{,}6 \times (4{,}5 + 5{,}5) = 3{,}6 \times 10 = 36$.
\item Đúng. $Y = 40 - 4 = 36$.
\item Đúng. Vì $X = 36$ và $Y = 36$ nên $X = Y$.
\item Đúng. $36 \times 36 = 1\,296$.
\end{itemchoice}
}
\end{ex}

\begin{ex} %[5D6V31]
Một gia đình thu hoạch hai thửa ruộng: thửa thứ nhất thu được $1{,}8$ tấn lúa; thửa thứ hai thu được nhiều hơn thửa thứ nhất $400\text{ kg}$ lúa.
\choiceTF
{\True Đổi $1{,}8\text{ tấn} = 1\,800\text{ kg}$}
{\True Thửa ruộng thứ hai thu hoạch được $2\,200\text{ kg}$ lúa}
{\True Cả hai thửa ruộng thu hoạch được tất cả $4$ tấn lúa}
{Nếu bán đi $60\%$ tổng số lúa với giá $8\,500$ đồng/kg thì thu được số tiền là $25\,000\,000$ đồng}
\loigiai{
\begin{itemchoice}
\item Đúng. $1{,}8\text{ tấn} = 1\,800\text{ kg}$.
\item Đúng. Thửa hai: $1\,800 + 400 = 2\,200\text{ (kg)}$.
\item Đúng. Cả hai thửa: $1\,800 + 2\,200 = 4\,000\text{ (kg)} = 4\text{ tấn}$.
\item Sai. Số lúa bán là: $4\,000 \times 60\% = 2\,400\text{ (kg)}$. Số tiền thu được là: $2\,400 \times 8\,500 = 20\,400\,000\text{ (đồng)} 
eq 25\,000\,000\text{ đồng}$.
\end{itemchoice}
}
\end{ex}

\textbf{Phần 3. Trắc nghiệm trả lời ngắn}
\setcounterpageref{ex}{0}

\begin{ex} %[5D6V31]
Tính giá trị của biểu thức: $1{,}25 \times 32 \times 2{,}5$.
\shortans{100}
\loigiai{
Tách $32 = 8 \times 4$:
$$(1{,}25 \times 8) \times (4 \times 2{,}5) = 10 \times 10 = 100.$$
}
\end{ex}

\begin{ex} %[5D6V31]
Tìm số tự nhiên $x$ thỏa mãn: $2{,}4 \times 5 < x < 3{,}5 \times 4$.
\shortans{13}
\loigiai{
Ta có: $2{,}4 \times 5 = 12$ và $3{,}5 \times 4 = 14$.\\
Do đó: $12 < x < 14$.\\
Vì $x$ là số tự nhiên nên $x = 13$.
}
\end{ex}

\begin{ex} %[5D6V31]
May 6 bộ rèm cửa hết $28{,}8\text{ m}$ vải. Hỏi may 15 bộ rèm cửa cùng loại thì cần tất cả bao nhiêu mét vải?
\shortans{72}
\loigiai{
May một bộ rèm cửa hết số mét vải là:
$28{,}8 : 6 = 4{,}8\text{ (m)}$.\\
May 15 bộ rèm cửa hết số mét vải là:
$4,8 \times 15 = 72\text{ (m)}$.
}
\end{ex}

\begin{ex} %[5D6V31]
Một ô tô trong 2 giờ đầu đi được $85{,}5\text{ km}$, trong 3 giờ sau đi được $139{,}5\text{ km}$. Hỏi trung bình mỗi giờ ô tô đó đi được bao nhiêu ki-lô-mét?
\shortans{45}
\loigiai{
Tổng quãng đường ô tô đi được là:
$85{,}5 + 139{,}5 = 225\text{ (km)}$.\\
Tổng thời gian đi là:
$2 + 3 = 5\text{ (giờ)}$.\\
Trung bình mỗi giờ ô tô đi được:
$225 : 5 = 45\text{ (km)}$.
}
\end{ex}

\begin{ex} %[5D6V31]
Tính nhanh giá trị của biểu thức: $45{,}6 : 0{,}25 + 45{,}6 \times 96$.
\shortans{4560}
\loigiai{
Vì chia cho $0{,}25$ là nhân với $4$ nên:
$$45{,}6 \times 4 + 45{,}6 \times 96 = 45{,}6 \times (4 + 96) = 45{,}6 \times 100 = 4\,560.$$
}
\end{ex}

\begin{ex} %[5D6V31]
Hai thùng dầu chứa tất cả $80\text{ l}$ dầu. Nếu chuyển $5\text{ l}$ dầu từ thùng thứ nhất sang thùng thứ hai thì lượng dầu ở hai thùng bằng nhau. Hỏi lúc đầu thùng thứ nhất chứa bao nhiêu lít dầu?
\shortans{45}
\loigiai{
Lúc đầu thùng thứ nhất chứa nhiều hơn thùng thứ hai số lít dầu là:
$5 \times 2 = 10\text{ (l)}$.\\
Lúc đầu thùng thứ nhất chứa số lít dầu là:
$(80 + 10) : 2 = 45\text{ (l)}$.
}
\end{ex}
